\SourceLeaf{381}%
auf den Parallelkreis, in welchem der Krümmungshalbmesser des Meridians dem Halbmesser eines Kreises gleich ist, dessen Umfang dem ganzen elliptischen Meridian gleich wird: die Breite dieses Parallelkreises findet sich, für die zum Grunde gelegte Abplattung 0,00324, \(45^0\,3'\,29''2\). Die geodaetischen Messungen ergeben den ganzen Bogen dieses Parallelkreises zwischen den Meridianen von Mailand (Sternwarte) und von Usson zu 475121,06 Meter, während die Beobachtungen der Pulversignale für den Längenunterschied \(6^0\,1'\,41''7\) gegeben haben. Man kann diese Zahlen als das Hauptresultat der Messungen betrachten. Die Vergleichung eines solchen Längengradbogens mit dem Resultat einer Breitengradmessung kann, theoretisch genommen, die Bestimmung der Erdabplattung geben: die Verf. finden aus einer solchen Vergleichung ihres Resultats mit dem Bogen von Greenwich bis Formentera die Abplattung \(\frac{1}{241}\). Wie wenig Zuverlässigkeit aber auf diese Weise erreicht werden kann, zeigt sich am auffallendsten, wenn man anstatt des ganzen Bogens die einzelnen Stücke auf ähnliche Art behandelt. Ref. findet so aus der Vergleichung desselben Meridianbogens mit dem Stück d’Usson --- Colombier die Abplattung \(\frac{1}{169}\), mit dem zweiten Stück Colombier --- Mont Cenis \(\frac{1}{226}\), mit dem vierten Turin --- Mailand \(\frac{1}{110}\), mit dem dritten Stück Mont Cenis --- Turin hingegen eine Allongation \(\frac{1}{82}\). In dieser Beziehung ist also hiervon für die schärfere Bestimmung der Erddimensionen wenig zu erwarten: allein desto wichtiger sind die Resultate, indem sie eine neue Bestätigung der Unregelmässigkeit der Erdfigur liefern, die sich gerade in Oberitalien im grössten Massstabe zeigt. Am deutlichsten treten diese Unregelmässigkeiten hervor, wenn man die astronomisch bestimmten Längenunterschiede mit den aus den geodaetischen Messungen, nach einer plausibeln Hypothese über die Erdfigur im Grossen, berechneten vergleicht. Die Verf. haben diese Rechnung mit der Abplattung 0,00324 und dem Aequatorshalbmesser 6376986 Meter geführt: auf diese Weise ergeben sich die westlichen Längenunterschiede mit Mailand in Zeit
\begin{center}
\begin{tabular}{l|r@{\quad}l|r@{\quad}l|r}
&\multicolumn{2}{c|}{astronomisch}&\multicolumn{2}{c|}{geodaetisch}&Unterschied\\
Turin & \(5'\) & \(58''85\) & \(6'\) & \(0''93\) & \(-2''08\)\\
Mont Cenis & \(9\) & \(0,20\) & \(8\) & \(59,49\) & \(+0,71\)\\
Colombier & \(13\) & \(44,23\) & \(13\) & \(43,84\) & \(+0,39\)\\
D’Usson & \(24\) & \(6,78\) & \(24\) & \(8,02\) & \(-1,24\)
\end{tabular}
\end{center}
